% Part: computability
% Chapter: recursive-functions
% Section: halting-problem

\documentclass[../../../include/open-logic-section]{subfiles}

\begin{document}

\olfileid[tr]{cmp}{rec}{hlt}
\olsection{Durma Problemi}

Genel olarak \emph{durma problemi}, hesaplanabilir bir fonksiyonun belirtimi
$e$ (örneğin programı) ile bir $n$ sayısı verildiğinde, bu fonksiyonun $n$
girdisindeki hesabının durup durmadığına, yani bir sonuç üretip üretmediğine
karar verme problemidir. Alan Turing, bu problemin kendisinin hesaplanabilir
bir fonksiyonla çözülemeyeceğini ünlü biçimde ispatladı. Yani
\[
  h(e,n)=\begin{cases}
    1 & \text{$e$ hesabı $n$ girdisinde duruyorsa}\\
    0 & \text{aksi durumda}
  \end{cases}
\]
fonksiyonu hesaplanabilir değildir.

Kısmi özyinelemeli fonksiyonlar bağlamında bir program belirtiminin rolünü,
Kleene normal biçim teoreminde verilen $e$ indisi oynayabilir. $f$ kısmi
özyinelemeli bir fonksiyonsa normal biçim teoremindeki denklemin geçerli
olduğu her $e$, $f$'nin bir indisidir. Bir $e$ sayısı verildiğinde normal
biçim teoremi
\[
  \cfind{e}(x)\simeq U(\mu s\;T(e,x,s))
\]
fonksiyonunun kısmi özyinelemeli olduğunu ve her kısmi özyinelemeli
$f\colon\Nat\to\Nat$ için bütün $x\in\Nat$ değerlerinde
$\cfind{e}(x)\simeq f(x)$ olacak bir $e\in\Nat$ bulunduğunu söyler. Aslında
her böyle $f$ için yalnızca bir değil, sonsuz sayıda böyle $e$ vardır.
\emph{Durma fonksiyonu}~$h$ şöyle tanımlanır:
\[
  h(e,x)=\begin{cases}
    1 & \text{$\cfind{e}(x)\fdefined$ ise}\\
    0 & \text{aksi durumda.}
  \end{cases}
\]
$\cfind{e}(x)\fundefined$ ise $h(e,x)=0$ olduğuna; kaynağın indis
gösterimini kısmi bıraktığı okumada $e$ hiçbir kısmi özyinelemeli fonksiyonun
indisi değilse de aynı değerin verildiğine dikkat edin.

\begin{thm}
\ollabel{thm:halting-problem}
Durma fonksiyonu $h$ kısmi özyinelemeli değildir.
\end{thm}

\begin{proof}
$h$ kısmi özyinelemeli olsaydı
\[
  d(y)=\begin{cases}
    1 & \text{$h(y,y)=0$ ise}\\
    \umin{x}{x\neq x} & \text{aksi durumda}
  \end{cases}
\]
fonksiyonunu tanımlayabilirdik. Hiçbir $x$ sayısı $x\neq x$ koşulunu
sağlamadığından $\umin{x}{x\neq x}$ yoktur ve
$d(y)\fundefined$ ancak ve ancak $h(y,y)\neq0$ olduğunda geçerlidir.
Tanımdan şunlar çıkar:
\begin{enumerate}
\item $d(y)\fdefined$ ancak ve ancak $\cfind{y}(y)\fundefined$ olduğunda ya
  da $y$ hiçbir kısmi özyinelemeli fonksiyonun indisi olmadığında geçerlidir.
\item $d(y)\fundefined$ ancak ve ancak $\cfind{y}(y)\fdefined$ olduğunda
  geçerlidir.
\end{enumerate}
$h$ kısmi özyinelemeli olsaydı $d$ de kısmi özyinelemeli olurdu. O zaman
Kleene normal biçim teoremi uyarınca $d$'nin bir $e_d$ indisi bulunurdu.
$h(e_d,e_d)$ değerini ele alalım. İki olası durum vardır: $0$ ve $1$.
\begin{enumerate}
\item $h(e_d,e_d)=1$ ise $\cfind{e_d}(e_d)\fdefined$ olur. Fakat
  $\cfind{e_d}\simeq d$ ve $d(e_d)$ ancak ve ancak $h(e_d,e_d)=0$ ise
  tanımlıdır. Dolayısıyla $h(e_d,e_d)\neq1$.
\item $h(e_d,e_d)=0$ ise ya $e_d$ hiçbir kısmi özyinelemeli fonksiyonun
  indisi değildir ya da bir indistir ve $\cfind{e_d}(e_d)\fundefined$ olur.
  Fakat yine $\cfind{e_d}\simeq d$ ve $d(e_d)$ ancak ve ancak
  $\cfind{e_d}(e_d)\fdefined$ ise tanımsızdır.
\end{enumerate}
Sonuçta $e_d$ gerçekte bir kısmi özyinelemeli fonksiyonun indisi olamaz.
Oysa $h$ kısmi özyinelemeli olsaydı $d$ de öyle olur ve $e_d$'yi onun bir
indisi olarak tanımlamamız geçerli olurdu. Öyleyse $h$ kısmi özyinelemeli
olamaz.
\end{proof}

\end{document}

